\documentclass[11pt,a4paper]{article}
\usepackage[T1]{fontenc}
\usepackage{lmodern}
\usepackage[a4paper,margin=1in]{geometry}
\usepackage{amsmath,amssymb,amsthm}
\usepackage{alltt}
\usepackage{booktabs}
\usepackage{array}
\usepackage[hidelinks,pdftitle={A proof of Conjecture 213 of Written on the Wall II},pdfauthor={Jiri Chmel}]{hyperref}
\usepackage{xurl}

\newtheorem{theorem}{Theorem}
\newtheorem{lemma}[theorem]{Lemma}
\newtheorem{proposition}[theorem]{Proposition}
\newtheorem{conjecture}{Conjecture}
\renewcommand{\theconjecture}{213}
\theoremstyle{remark}
\newtheorem{remark}[theorem]{Remark}

\newcommand{\lm}{\operatorname{lm}}
\newcommand{\Gc}{\overline{G}}
\newcommand{\Fc}{\overline{F}}
\newcommand{\ceil}[1]{\lceil #1\rceil}
\newcommand{\flr}[1]{\lfloor #1\rfloor}

\title{A proof of Conjecture 213 of Written on the Wall II}
\author{Ji\v{r}\'{\i} Chmel}
\date{22 November 2026\\[1ex]
\small Preprint, version 1. It has not yet been fully reviewed by the author;\\
a version reviewed by the author will follow.}

\begin{document}
\maketitle

\begin{abstract}
Conjecture 213 of Written on the Wall II, a list of conjectures made by the program Graffiti.pc,
states that a connected graph $G$ on $n\ge 2$ vertices has a Hamiltonian path if the
2-domination number of its complement exceeds the lower median of the degree sequence of the
complement by at least two. We prove it. The core of the proof is the inequality
$\gamma_2(\Gc)\le\lm(\Gc)+1$ for every connected, non-complete graph $G$ that is closed in the
sense of Bondy and Chv\'atal, that is, $d(u)+d(v)\le n-2$ for all non-adjacent vertices $u$ and
$v$. The conjecture follows because adding an edge to $G$ cannot decrease
$\gamma_2(\Gc)-\lm(\Gc)$, and an edge-maximal graph without a Hamiltonian path is closed. The
inequality is proved by an explicit polynomial-time construction based on a selection lemma for
set systems. It is best possible, and so is the constant~2 in the conjecture. The proof has been
formalised in Lean~4 with Mathlib and checked by the Lean kernel.
\end{abstract}

\section{Introduction}

Written on the Wall II~\cite{wowii} is E.~DeLaVi\~na's list of conjectures made by her program
Graffiti.pc~\cite{gpc}. Its entry 213, dated 12 January 2006 and marked open in the most recent version we could consult, concerns
two invariants of the complement $\Gc$ of a graph $G$. The first is the 2-domination number
$\gamma_2(\Gc)$: the least size of a set $D$ of vertices such that every vertex outside $D$ has
at least two neighbours in~$D$~\cite{fj,dlpw}. The second is the lower median $\lm(\Gc)$ of the
degree sequence of $\Gc$, the $\ceil{n/2}$-th smallest degree. Section~\ref{sec:defs} says why
these are the definitions of the catalogue.

\begin{conjecture}[\cite{wowii}]
If $G$ is a simple connected graph with $n>1$ vertices such that
$2+\lm(\Gc)\le\gamma_2(\Gc)$, then $G$ has a Hamiltonian path.
\end{conjecture}

We prove it.

\begin{theorem}\label{thm:main}
Every connected graph $G$ on $n\ge 2$ vertices with $\gamma_2(\Gc)\ge\lm(\Gc)+2$ has a
Hamiltonian path.
\end{theorem}

The hypothesis holds for many graphs. The complete graph $K_n$ with $n\ge 2$ satisfies it, since
$\gamma_2(\overline{K_n})=n$ and $\lm(\overline{K_n})=0$. Counted up to isomorphism, the
connected graphs on $n=2,3,\dots,10$ vertices that satisfy it number 1, 2, 5, 11, 57, 261, 2248,
27724 and 874505 (Section~\ref{sec:verify}).

Call a graph $G$ on $n$ vertices \emph{closed} if $d(u)+d(v)\le n-2$ for any two non-adjacent
vertices $u$ and $v$. In the language of Bondy and Chv\'atal~\cite{bc}, $G$ equals its own
$(n-1)$-closure. Theorem~\ref{thm:main} rests on the following inequality.

\begin{theorem}\label{thm:key}
Let $G$ be a closed, connected graph on $n\ge 2$ vertices that is not complete. Then
$\gamma_2(\Gc)\le\lm(\Gc)+1$. Moreover, a 2-dominating set of $\Gc$ with at most $\lm(\Gc)+1$
vertices can be found in time polynomial in~$n$.
\end{theorem}

Theorem~\ref{thm:main} follows from Theorem~\ref{thm:key} by a standard reduction
(Section~\ref{sec:reduction}). Adding an edge to $G$ cannot decrease $\gamma_2(\Gc)-\lm(\Gc)$.
A graph without a Hamiltonian path to which no edge can be added without creating one is closed,
by the closure argument of Bondy and Chv\'atal.

The proof of Theorem~\ref{thm:key} (Sections~\ref{sec:median}--\ref{sec:proofkey}) splits
according to $m=\lm(\Gc)$. If $2m\ge n$, a counting argument gives a 2-dominating set of size
$m+1$ in every graph. If $2m\le n-1$, closedness forces the vertices of degree at least $m$ to
form a clique, and the vertices of smaller degree to be adjacent to few of its non-universal
vertices. A selection lemma for set systems (Lemma~\ref{lem:select}) then picks the set.

Both constants are best possible (Section~\ref{sec:sharp}). For $s\ge 1$ and $n\ge 2s+2$ the
graph $F(n,s)=K_s\vee\bigl((s+1)K_1\cup K_{n-2s-1}\bigr)$ is closed, connected and not complete,
it has no Hamiltonian path, and $\gamma_2(\Fc)=\lm(\Fc)+1$. So the bound of
Theorem~\ref{thm:key} is attained, and the constant~2 in Conjecture~213 cannot be replaced
by~1.

The proof of Theorem~\ref{thm:main} has been formalised in Lean~4 with Mathlib and checked by
the Lean kernel (Section~\ref{sec:verify}). Appendix~\ref{app:lean} gives the formal statement.

\emph{Earlier work.} In October 2026 an agent of the AI Village project (Claude Opus 5)
studied Conjecture~213~\cite{aiv}. Its note restates the median condition in terms of the large
degrees of~$G$, proves several structural lemmas, checks the conjecture on every connected graph
with at most 10 vertices, and finds the graphs $F(2t+1,t-1)=K_{t-1}\vee(K_2\cup tK_1)$, $t\ge 2$, as a
family of near misses. It states that the conjecture is not proved for any order. Qing~\cite{qing}
proved Graffiti.pc's Conjecture 387, $\gamma_2(X)\le n-\operatorname{um}(X)+1$ for every nonempty graph $X$
on $n$ vertices, where $\operatorname{um}(X)$ is the upper median of the degree sequence; it
implies Lemma~\ref{lem:J} below, which, apart from a trivial case, is also a special case of an
older inequality of Pepper~\cite{pepper,hp}. Li~\cite{li} related the domination number of the complement to
traceability. We found no other
work on this conjecture. The end of the paper lists the sources we consulted.

\subsection{Use of generative AI}\label{sec:ai}

The proof was found, and this paper was written, by Claude Opus~5.5, a generative AI system made
by Anthropic, on 20 and 21 November 2026, in sessions directed by the author. Claude was used
through Claude Code, Anthropic's agentic coding tool. The author chose the
problem, set the working rules and decided which lines of work to pursue. Claude searched for
counterexamples, found the proof of Theorem~\ref{thm:key} and the reduction of
Theorem~\ref{thm:main} to it, wrote the programs and the Lean formalisation described in
Section~\ref{sec:verify}, and drafted the text. Further Claude instances, writing programs of their
own, refereed the proof by hand, re-checked its finite consequences, audited the formal statement
against the catalogue's definitions, and checked the references against primary sources. None of
this is external peer review. The Lean kernel's check of the formal proof does not depend on any AI
system; that the formal statement expresses Conjecture~213 was checked as described in
Section~\ref{sec:verify} and Appendix~\ref{app:lean}. Claude is not an author. The author takes
full responsibility for the content of this paper. This first version is a preprint, posted on
22 November 2026 before the author's own full review of the proofs; a version reviewed by the
author will replace it.

\section{Definitions and the reading of the statement}\label{sec:defs}

Graphs are finite and simple. For a graph $G=(V,E)$ with $n=|V|$ vertices, $N(v)$ is the
neighbourhood of a vertex $v$ and $d(v)=|N(v)|$ its degree. Two vertices are
\emph{non-adjacent} if they are distinct and not adjacent. The complement $\Gc$ has the same
vertex set, and two distinct vertices are adjacent in $\Gc$ if and only if they are non-adjacent
in $G$. We write $H=\Gc$ and $d_H(v)=n-1-d(v)$. A vertex of degree $n-1$ is \emph{universal}.

\emph{2-domination.} A set $D$ of vertices of a graph $X$ is \emph{2-dominating} if every vertex
of $X$ outside $D$ has at least two neighbours in $D$. The \emph{2-domination number}
$\gamma_2(X)$ is the least size of a 2-dominating set. The neighbours of $v$ in $H$ are the
vertices non-adjacent to $v$ in $G$. Hence
\begin{equation}\label{eq:twodom}
  D\text{ is 2-dominating in }H \iff
  \text{every }v\notin D\text{ is non-adjacent in }G\text{ to at least two vertices of }D.
\end{equation}
This is the case $k=2$ of $k$-domination~\cite{fj}. Definition 90 of the catalogue's definitions
file requires each vertex to be in $D_2$ or adjacent to 2 vertices of~$D_2$~\cite{wowii}. The
paper of DeLaVi\~na, Larson, Pepper and Waller on Graffiti.pc's conjectures about
2-domination~\cite{dlpw} defines $k$-dominating sets with ``at least $k$'' neighbours. The
distinction matters, because the conjecture would be false if every vertex outside $D$ had to
have \emph{exactly} two neighbours in $D$. For example, let $G$ have vertex set $\{0,\dots,6\}$
and edges $06,16,25,26,34,35,45,46,56$. It is connected and has no Hamiltonian path, because the
vertices 0 and 1 are leaves at the same vertex. In $H$ the neighbourhoods are
$N_H(0)=\{1,2,3,4,5\}$, $N_H(1)=\{0,2,3,4,5\}$, $N_H(2)=\{0,1,3,4\}$, $N_H(3)=\{0,1,2,6\}$,
$N_H(4)=\{0,1,2\}$, $N_H(5)=\{0,1\}$ and $N_H(6)=\{3\}$. The degrees of $H$ are
$1,2,3,4,4,5,5$, so the lower median, the upper median and the term $h_5$ discussed below all
equal 4. The set $\{0,1,6\}$ is 2-dominating in $H$, so $\gamma_2(H)\le 3$, and $G$ does not
satisfy the hypothesis. Now let $D$ be a set in which every vertex outside $D$ has exactly two
$H$-neighbours in $D$. Then $6\in D$. If $5\notin D$, then $0,1\in D$, which forces first $3$,
then $2$, then $4$ into $D$, so $|D|=6$. If $5\in D$ and neither 0 nor 1 is in $D$, then
vertex 0 has exactly one of $2,3,4$ in $D$, and then vertex 4 or vertex 2 has only one
$H$-neighbour in $D$, which is impossible. If $5\in D$ and, say, $0\in D$ (0 and 1 play
symmetric roles), then $1\in D$, since otherwise $2,3,4\notin D$ and vertex 4 would have only the
$H$-neighbour 0 in $D$; then $3$, $2$ and $4$ are forced into $D$ as before, and $D=V$. So such a
set has at least $6=4+2$ vertices, and under the ``exactly two'' reading $G$ would satisfy the
hypothesis, whichever median is used.

\emph{Lower median.} Let $h_1\le h_2\le\dots\le h_n$ be the degrees of $H$ in nondecreasing
order. We use
\[
  \lm(H)=h_{\ceil{n/2}},
\]
which is $h_{n/2}$ for even $n$ and the median $h_{(n+1)/2}$ for odd $n$. For even $n$ this is
exactly the catalogue's definition 68. For odd $n$ the definition says that the lower and
the upper median both equal the median. When it describes the median of an odd sequence, however,
it names the term $d_{(n+1)/2+1}$. This is a slip: the median of $d_1\le\dots\le d_n$ for odd $n$
is $d_{(n+1)/2}$, and~\cite{dlpw} defines the upper median through the statistical median. The slip does not affect our
result. The upper median $h_{\flr{n/2}+1}$ and the term $h_{(n+3)/2}$ named by the slip are both
at least $h_{\ceil{n/2}}$. Under either reading the hypothesis of Conjecture~213 is therefore
stronger, and Theorem~\ref{thm:main} still applies.

A \emph{Hamiltonian path} is a path through all vertices, and a graph is \emph{traceable} if it
has one. A Hamiltonian cycle is a cycle through all vertices (of a graph with at least three
vertices).

\emph{Notation for Sections~\ref{sec:median}--\ref{sec:proofkey}.} $G$ is a graph on $n\ge 2$
vertices, $H=\Gc$ and $m=\lm(H)$. $U$ is the set of universal vertices of $G$, and $u=|U|$.

\section{Counting with the median}\label{sec:median}

\begin{lemma}\label{lem:median}
Let $t$ be an integer.
\begin{itemize}
\item[(i)] If $m\le t$, then at least $\ceil{n/2}$ vertices have $d_H\le t$.
\item[(ii)] If $m\ge t$, then at least $\flr{n/2}+1$ vertices have $d_H\ge t$.
\end{itemize}
\end{lemma}

\begin{proof}
(i) $h_1\le\dots\le h_{\ceil{n/2}}=m\le t$. (ii) $t\le m=h_{\ceil{n/2}}\le\dots\le h_n$, and
there are $n-\ceil{n/2}+1=\flr{n/2}+1$ of these terms. In both cases the terms are the degrees
of distinct vertices.
\end{proof}

\section{The case \texorpdfstring{$2m\ge n$}{2m >= n}}\label{sec:caseJ}

\begin{lemma}\label{lem:J}
Every graph $G$ on $n\ge 2$ vertices with $2m\ge n$ satisfies $\gamma_2(H)\le m+1$.
\end{lemma}

\begin{proof}
Let $W=\{v: d_H(v)\ge m\}$. By Lemma~\ref{lem:median}(ii) with $t=m$,
$|W|\ge\flr{n/2}+1\ge 1$, so some vertex has $d_H\ge m$, and $m\le n-1$. Put $a=n-m-1\ge 0$. Since
$m\ge n/2$, $a\le n/2-1<\flr{n/2}+1\le|W|$. Choose $A\subseteq W$ with $|A|=a$ and let
$D=V\setminus A$, so that $|D|=m+1$. A vertex $x\in A$ has at least $m$ neighbours in $H$, of
which at most $|A|-1=n-m-2$ lie in $A$. So at least $m-(n-m-2)=2m-n+2\ge 2$ of them lie in $D$.
Every vertex outside $D$ lies in $A$, so $D$ is 2-dominating in~$H$.
\end{proof}

For $m\le n-2$, Lemma~\ref{lem:J} is the case $k=2$, $j=n-m-1\ge 1$ of Pepper's inequality
$\gamma_k(X)+\alpha_j(X_{\ge k+j-1})\le |V(X)|$~\cite{pepper} (see \cite[Theorem 3(i)]{hp} and
\cite[Theorem 2.29]{hv}).
Here $X_{\ge t}$ is the subgraph induced by the vertices of degree at least $t$, and $\alpha_j$ is the
largest order of a vertex set inducing a subgraph of maximum degree less than $j$. With $X=H$ we
have $k+j-1=n-m\le m$, so $W$ lies in $H_{\ge n-m}$, and any $j$ vertices of $W$ form such a set. It also follows from Qing's theorem~\cite{qing} applied to $X=H$: the upper
median of the degree sequence of $H$ is at least $m$, and $n-m+1\le m+1$ when $2m\ge n$. For
$m=n-1$ the lemma is trivial, since $\gamma_2(H)\le n$.

\section{Closed graphs with \texorpdfstring{$2m\le n-1$}{2m <= n-1}}\label{sec:structure}

In this section $G$ is closed, connected and not complete, and $2m\le n-1$. Let
\[
  T=\{v: d(v)\ge n-1-m\},\quad Q=\{v: d(v)\ge m\},\quad Y=V\setminus Q,\quad P=Q\setminus U,
  \quad r=m-u+1 .
\]
Note that $T=\{v: d_H(v)\le m\}$.

\begin{lemma}\label{lem:structure}
In this setting the following hold.
\begin{itemize}
\item[(a)] $|T|\ge\ceil{n/2}$ and $T\subseteq Q$.
\item[(b)] $Q$ is a clique of $G$.
\item[(c)] $Y\neq\emptyset$ and $U\subseteq Q$. Every vertex $y\in Y$ is adjacent to every vertex
  of $U$, and $u\le d(y)\le m-1$. In particular $u\le m-1$ and $r\ge 2$.
\item[(d)] Every non-neighbour of a vertex of $Q$ lies in $Y$. Every vertex of $P$ has a
  non-neighbour.
\item[(e)] Every vertex of $P$ has at least two non-neighbours.
\item[(f)] Every vertex of $Y$ has at most $r-2$ neighbours in $P$.
\item[(g)] $V$ is the disjoint union of $U$, $P$ and $Y$.
\item[(h)] If $|D|\ge m+1$ and $y\in Y\setminus D$, then $y$ is non-adjacent to at least two
  vertices of~$D$.
\end{itemize}
\end{lemma}

\begin{proof}
(a) By Lemma~\ref{lem:median}(i) with $t=m$, at least $\ceil{n/2}$ vertices have $d_H\le m$,
that is, $d\ge n-1-m$; these form $T$. If $v\in T$, then $d(v)\ge n-1-m\ge m$ because
$2m\le n-1$. So $T\subseteq Q$.

(b) Suppose that $a,b\in Q$ are non-adjacent. Closedness gives $d(a)+d(b)\le n-2$. If $a\in T$,
then $d(a)+d(b)\ge(n-1-m)+m=n-1$, a contradiction; the same holds if $b\in T$. So
$a,b\notin T$. Every $t\in T$ differs from $a$ and lies in $Q$ by (a), and it is adjacent to
$a$, since otherwise the pair $t,a$ would give the contradiction just obtained. Hence
$d(a)\ge|T|\ge\ceil{n/2}$, and likewise $d(b)\ge\ceil{n/2}$. Then $d(a)+d(b)\ge n>n-2$, a
contradiction.

(c) If $Y$ were empty, then $Q=V$ would be a clique by (b), and $G$ would be complete. A
universal vertex has degree $n-1\ge m$, so $U\subseteq Q$. Let $y\in Y$. Then
$d(y)\le m-1<n-1$, so $y$ is not universal, and every universal vertex is adjacent to $y$. So
$U\subseteq N(y)$, and $u\le d(y)\le m-1$. As $Y\ne\emptyset$, $u\le m-1$ and $r\ge 2$.

(d) If $q\in Q$ and $w$ is a non-neighbour of $q$, then $w\notin Q$ by (b), so $w\in Y$. A vertex
of $P$ is not universal, so it has a non-neighbour.

(e) Let $q\in P$, and let $y$ be a non-neighbour of $q$, which exists by (d). Closedness gives
$d(q)+d(y)\le n-2$. As $G$ is connected and $n\ge 2$, $d(y)\ge 1$. Hence $d(q)\le n-3$, and $q$
has $n-1-d(q)\ge 2$ non-neighbours.

(f) Let $y\in Y$. By (c), $U\subseteq N(y)$, and $U\cap P=\emptyset$. So
$|N(y)\cap P|\le d(y)-u\le m-1-u=r-2$.

(g) $U\subseteq Q$ by (c), $P=Q\setminus U$ and $Y=V\setminus Q$.

(h) $y$ has $d(y)\le m-1$ neighbours, so at most $m-1$ vertices of $D$ are adjacent to $y$.
The remaining $|D|-(m-1)\ge 2$ vertices of $D$ differ from $y$, because $y\notin D$, and are not
adjacent to~$y$.
\end{proof}

\section{A selection lemma}\label{sec:select}

\begin{lemma}\label{lem:select}
Let $r\ge 2$ be an integer, let $P$ be a finite set, let $Y'$ be a finite set with $|Y'|\ge r$,
and for each $y\in Y'$ let $A_y\subseteq P$ with $|A_y|\le r-2$. Then there are $B\subseteq P$
and $R\subseteq Y'$ with $|B|+|R|=r$ and $|R|\ge 2$ such that
\[
  |\{y\in R: q\in A_y\}|\le|R|-2\qquad\text{for every }q\in P\setminus B .
\]
\end{lemma}

\begin{proof}
Induction on $r$. If $r=2$, every $A_y$ is empty. Take $B=\emptyset$ and any $R\subseteq Y'$
with $|R|=2$: then $|\{y\in R: q\in A_y\}|=0=|R|-2$ for every $q\in P$.

Let $r\ge 3$, and assume the statement for $r-1$. Choose $R_0\subseteq Y'$ with $|R_0|=r$. If
$|\{y\in R_0: q\in A_y\}|\le r-2$ for every $q\in P$, take $B=\emptyset$ and $R=R_0$.
Otherwise there is $x\in P$ with $|\{y\in R_0: x\in A_y\}|\ge r-1$. Let
$Y''=\{y\in R_0: x\in A_y\}$, so that $|Y''|\ge r-1$. Put $P'=P\setminus\{x\}$ and
$A'_y=A_y\setminus\{x\}$ for $y\in Y''$. Then $A'_y\subseteq P'$ and
$|A'_y|=|A_y|-1\le(r-1)-2$. Since $r-1\ge 2$, the induction hypothesis applies to
$r-1$, $P'$, $Y''$ and the sets $A'_y$. It gives $B'\subseteq P'$ and $R\subseteq Y''$ with
$|B'|+|R|=r-1$ and $|R|\ge 2$, such that $|\{y\in R: q\in A'_y\}|\le|R|-2$ for every
$q\in P'\setminus B'$. Put $B=B'\cup\{x\}$. As $x\notin P'$, $|B|+|R|=r$. Let
$q\in P\setminus B$. Then $q\neq x$ and $q\in P'\setminus B'$. For $y\in R$ we have $q\in A_y$ if
and only if $q\in A'_y$, because $q\ne x$. So $|\{y\in R: q\in A_y\}|\le|R|-2$.
\end{proof}

\begin{remark}\label{rem:procedure}
Unrolled, the proof is a procedure. Start with any $r$ indices $R$ and $B=\emptyset$. While some
$q\in P\setminus B$ has $|\{y\in R: q\in A_y\}|\ge|R|-1$, add $q$ to $B$ and replace $R$ by
$|R|-1$ of the indices $y\in R$ with $q\in A_y$. Throughout, $|B|+|R|=r$, and $B\subseteq A_y$
for every $y\in R$. A round starts with $|R|\ge 2$ and leaves $R$ non-empty, so afterwards
$|B|\le r-2$, because $|A_y|\le r-2$, and hence $|R|\ge 2$ again. So there are at most $r-2$
rounds, and the final pair $(B,R)$ satisfies Lemma~\ref{lem:select}. The set $B$ cannot always
be empty: if $r\ge 3$, $|Y'|=r$ and all the sets $A_y$ equal one set $S$ with $|S|=r-2$, then every
admissible choice has $B=S$ and $|R|=2$.
\end{remark}

\section{Proof of Theorem~\ref{thm:key}}\label{sec:proofkey}

Let $G$ be closed, connected and not complete, with $n\ge 2$ vertices, and let $m=\lm(H)$.

\emph{Case 1: $2m\ge n$.} Lemma~\ref{lem:J} gives $\gamma_2(H)\le m+1$.

\emph{Case 2: $2m\le n-1$.} Then Lemma~\ref{lem:structure} applies.

\emph{Case 2a: $|Y|\ge r$.} Apply Lemma~\ref{lem:select} with $r$, $P$, $Y'=Y$ and
$A_y=N(y)\cap P$. This is possible because $r\ge 2$ by Lemma~\ref{lem:structure}(c) and
$|A_y|\le r-2$ by Lemma~\ref{lem:structure}(f). It gives $B\subseteq P$ and $R\subseteq Y$ with
$|B|+|R|=r$. Let $D=U\cup B\cup R$. These three sets are disjoint by
Lemma~\ref{lem:structure}(g), so $|D|=u+r=m+1$. Let $v\notin D$. Then $v\notin U$, so $v\in Y$
or $v\in P$. If $v\in Y$, then $v$ is non-adjacent to two vertices of $D$ by
Lemma~\ref{lem:structure}(h). If $v\in P$, then $v\notin B$. For $y\in R$, $v\in A_y$ if and only
if $v$ is adjacent to $y$, so by Lemma~\ref{lem:select} $v$ is adjacent to at most $|R|-2$
vertices of $R$. So at least two
vertices of $R$ are non-adjacent to $v$; they lie in $D$. By~\eqref{eq:twodom}, $D$ is
2-dominating in $H$.

\emph{Case 2b: $|Y|\le r-1=m-u$.} Put $c=m+1-u-|Y|\ge 1$. By Lemma~\ref{lem:structure}(g),
$|P|=n-u-|Y|=c+(n-m-1)\ge c$. Choose $B_0\subseteq P$ with $|B_0|=c$ and let
$D=U\cup Y\cup B_0$, so that $|D|=m+1$. A vertex outside $D$ lies in $P$. By
Lemma~\ref{lem:structure}(e) it has at least two non-neighbours, and by
Lemma~\ref{lem:structure}(d) all of them lie in $Y\subseteq D$. By~\eqref{eq:twodom}, $D$ is
2-dominating in $H$.

In every case $\gamma_2(H)\le m+1$. Every step is effective: the degrees, $m$ and the sets $T$,
$Q$, $U$, $P$, $Y$ take polynomial time, the set $A$ of Lemma~\ref{lem:J} is found by sorting
the degrees, and the procedure of Remark~\ref{rem:procedure} has at most $r-2$ rounds of
polynomial cost. This proves Theorem~\ref{thm:key}. \qed

\section{From Theorem~\ref{thm:key} to Conjecture 213}\label{sec:reduction}

\begin{lemma}\label{lem:mono}
Let $e$ be a non-edge of a graph $G$. The complement of $G+e$ is $H-e$, and
$\gamma_2(H-e)\ge\gamma_2(H)$ and $\lm(H-e)\le\lm(H)$. In particular, if $G$ satisfies
$\gamma_2(\Gc)\ge\lm(\Gc)+2$, then so does $G+e$.
\end{lemma}

\begin{proof}
A 2-dominating set of $H-e$ is 2-dominating in $H$, which has more edges. Removing $e$ does not
increase any degree. If $x_v\le z_v$ for every vertex $v$, then for every $k$ the $k$-th
smallest $x_v$ is at most the $k$-th smallest $z_v$: the $k$ vertices with the $k$ smallest
values $z_v$ all have $x_v$ at most the $k$-th smallest $z_v$. So $\lm(H-e)\le\lm(H)$.
\end{proof}

\begin{lemma}[Ore~\cite{ore}; Bondy and Chv\'atal~\cite{bc}]\label{lem:ore}
Let $F$ be a graph on $N\ge 3$ vertices, and let $a,b$ be non-adjacent vertices with
$d_F(a)+d_F(b)\ge N$. If $F+ab$ has a Hamiltonian cycle, then so has $F$.
\end{lemma}

\begin{proof}
Suppose that $F+ab$ has a Hamiltonian cycle and $F$ has none. Then the cycle uses $ab$, and
$F$ has a Hamiltonian path $a=v_1,v_2,\dots,v_N=b$. Let
$S_a=\{i: 1\le i\le N-1,\ a\sim v_{i+1}\}$ and $S_b=\{i:1\le i\le N-1,\ v_i\sim b\}$. All
neighbours of $a$ lie among $v_2,\dots,v_N$ and all neighbours of $b$ among
$v_1,\dots,v_{N-1}$, so $|S_a|=d_F(a)$ and $|S_b|=d_F(b)$. Both sets lie in $\{1,\dots,N-1\}$,
and $|S_a|+|S_b|\ge N$, so they have a common element $i$. Then
$v_1v_2\cdots v_iv_Nv_{N-1}\cdots v_{i+1}v_1$ is a Hamiltonian cycle of $F$: it uses the edges
$v_iv_N=v_ib$ and $v_{i+1}v_1=v_{i+1}a$, and $i\ne 1,N-1$ because $a$ and $b$ are
non-adjacent. This is a contradiction.
\end{proof}

\begin{lemma}\label{lem:closure}
Let $G$ be a graph on $n\ge 2$ vertices that has no Hamiltonian path, such that $G+e$ has one for
every non-edge $e$. Then $G$ is closed.
\end{lemma}

\begin{proof}
Let $G^*$ be $G$ together with a new vertex $z$ adjacent to all vertices of $G$; it has
$N=n+1\ge 3$ vertices. $G$ has a Hamiltonian path if and only if $G^*$ has a Hamiltonian cycle:
a path $v_1\cdots v_n$ gives the cycle $zv_1\cdots v_nz$, and deleting $z$ from a Hamiltonian
cycle of $G^*$ leaves a Hamiltonian path of $G$. Suppose that $a,b$ are non-adjacent in $G$ and
$d(a)+d(b)\ge n-1$. In $G^*$ they are non-adjacent with degree sum $d(a)+d(b)+2\ge N$. As $G+ab$
has a Hamiltonian path, $(G+ab)^*=G^*+ab$ has a Hamiltonian cycle. By Lemma~\ref{lem:ore}, so
has $G^*$, and $G$ has a Hamiltonian path, a contradiction.
\end{proof}

Lemma~\ref{lem:closure} is the path form of the stability of the Hamiltonian property under the
$(n-1)$-closure, proved by the same construction in~\cite[Theorem 9.4]{bc}.

\begin{proof}[Proof of Theorem~\ref{thm:main}]
Let $G$ be connected with $n\ge 2$ vertices and $\gamma_2(\Gc)\ge\lm(\Gc)+2$, and suppose that
$G$ has no Hamiltonian path. Add non-edges to $G$ one at a time, each time only if the new graph
still has no Hamiltonian path, until no further non-edge can be added. The result is a graph
$G'$ on the same vertex set that contains $G$, has no Hamiltonian path, and gets one whenever a
non-edge is added. $G'$ is connected because it contains $G$, and it is not complete because
$K_n$ has a Hamiltonian path. By Lemma~\ref{lem:mono}, applied at each added edge,
$\gamma_2(\overline{G'})\ge\lm(\overline{G'})+2$. By Lemma~\ref{lem:closure}, $G'$ is closed.
This contradicts Theorem~\ref{thm:key}.
\end{proof}

Equivalently, for $n\ge 4$ one may pass to the $(n-1)$-closure of $G$~\cite{bc}, which has a
Hamiltonian path if and only if $G$ has one, and is closed.

\section{Sharpness and remarks}\label{sec:sharp}

\begin{proposition}\label{prop:sharp}
Let $s\ge 1$ and $n\ge 2s+2$. Let $S$, $I$ and $C$ be disjoint sets with $|S|=s$, $|I|=s+1$
and $|C|=n-2s-1$, and let $F=F(n,s)$ be the graph on $S\cup I\cup C$ in which $S\cup C$ is a
clique, every vertex of $S$ is adjacent to every vertex of $I$, and there are no other edges.
Then $F$ is closed, connected and not complete, it has no Hamiltonian path, and
$\gamma_2(\Fc)=s+2=\lm(\Fc)+1$.
\end{proposition}

\begin{proof}
The vertices of $S$ have degree $n-1$, those of $I$ degree $s$, and those of $C$ degree
$s+(n-2s-2)=n-s-2$. $F$ is connected because $s\ge 1$, and not complete because $|I|\ge 2$. The
non-adjacent pairs are the pairs inside $I$, with degree sum $2s\le n-2$, and the pairs $y\in I$,
$z\in C$, with degree sum $n-2$. So $F$ is closed. Deleting $S$ leaves $s+2$ components, the
$s+1$ vertices of $I$ and the clique $C\ne\emptyset$, whereas deleting $s$ vertices from a graph
with a Hamiltonian path leaves at most $s+1$ components. So $F$ has no Hamiltonian path.

In $\Fc$ the vertices of $S$ are isolated, $I$ is a clique, $C$ is independent, and every vertex
of $I$ is adjacent to every vertex of $C$. The degrees of $\Fc$ are $0$ ($s$ times), $s+1$
($n-2s-1$ times, on $C$) and $n-s-1$ ($s+1$ times, on $I$). As $n\ge 2s+2$, we have
$s+1\le n-s-1$ and $s+1\le\ceil{n/2}\le n-s-1$, so $\lm(\Fc)=h_{\ceil{n/2}}=s+1$. Every
2-dominating set $D$ of $\Fc$ contains $S$, whose vertices have no neighbours. If $D$ contained
at most one vertex of $I\cup C$, then some vertex of $I\cup C$, a set of $n-s\ge 3$ vertices,
would lie outside $D$ and have at most one neighbour in $D$. So $|D|\ge s+2$. Conversely, $S$
together with two vertices of $I$ is 2-dominating, since every other vertex of $I\cup C$ is
adjacent to both of them. So $\gamma_2(\Fc)=s+2$.
\end{proof}

The graphs $F(n,1)=K_1\vee(2K_1\cup K_{n-3})$, $n\ge 4$, show that the bound of
Theorem~\ref{thm:key} is attained for every $n\ge 4$, and that for every $n\ge 4$ the inequality
$1+\lm(\Gc)\le\gamma_2(\Gc)$ does not imply a Hamiltonian path. For $n\le 3$ every connected
graph has one. The family $F(2t+1,t-1)$, $t\ge 2$, is the near-miss family of~\cite{aiv}.

\begin{remark}
Closedness enters only in Lemma~\ref{lem:structure}(b) and~(e). Connectivity enters in
Lemma~\ref{lem:structure}(e), which Case~2b uses, and in the proof of Theorem~\ref{thm:main},
where $G'$ must be connected. Non-completeness enters in Lemma~\ref{lem:structure}(c) and in the
same proof. In the proof of Theorem~\ref{thm:key} the median enters only through
Lemma~\ref{lem:median}, and Lemma~\ref{lem:J} holds for all graphs; the reduction also uses that $\lm$
cannot increase when degrees decrease (Lemma~\ref{lem:mono}). Connectivity cannot be dropped from Theorem~\ref{thm:key}. For $n\ge 4$ the
graph $K_{n-1}\cup K_1$ is closed and not complete, its complement is the star $K_{1,n-1}$, and
$\gamma_2(K_{1,n-1})=n-1$, while $\lm(K_{1,n-1})=1$.
\end{remark}

\section{Verification}\label{sec:verify}

\emph{Formal proof.} Theorem~\ref{thm:main} is formalised in Lean~4~\cite{lean}, version
4.37.0-rc1, with Mathlib~\cite{mathlib} at revision \texttt{31ce06eaadf4}. The formalisation has
eight files with 1642 lines (namespace \texttt{W213}), and an audit file with 160 lines. The
theorem \texttt{W213.conjecture213Statement} proves the statement of Appendix~\ref{app:lean}.
For it and for \texttt{W213.conjecture213}, \texttt{\#print axioms} lists only
\texttt{propext}, \texttt{Classical.choice} and \texttt{Quot.sound}, the standard axioms of
Mathlib. No file uses \texttt{sorry}, \texttt{admit} or \texttt{native\_decide}, and none
declares an axiom. The formal proof follows this paper, with one difference: the closure
step (Lemmas~\ref{lem:ore} and~\ref{lem:closure}) is proved directly for Hamiltonian paths,
without the auxiliary vertex. Table~\ref{tab:lean} lists the correspondence. Neither the polynomial
running time in Theorem~\ref{thm:key} nor Proposition~\ref{prop:sharp} is formalised.

\begin{table}[ht]
\centering\small
\begin{tabular}{@{}>{\raggedright\arraybackslash}p{3.3cm}p{8.4cm}l@{}}
\toprule
This paper & Lean (namespace \texttt{W213}) & File\\
\midrule
Lemma~\ref{lem:median} & \texttt{lowerMedian\_le\_iff}, \texttt{le\_lowerMedian\_iff} & \texttt{Defs.lean}\\
Lemma~\ref{lem:J} & \texttt{lemmaJ}, \texttt{lemmaJ\_compl} & \texttt{CaseJ.lean}\\
Lemma~\ref{lem:structure} & \texttt{setT}, \texttt{setQ}, \texttt{setY}, \dots, and the lemmas about them & \texttt{Structure.lean}\\
Lemma~\ref{lem:select} & \texttt{selection\_lemma} & \texttt{Selection.lean}\\
Theorem~\ref{thm:key}, the bound & \texttt{theorem1}, \texttt{corollary2} & \texttt{Theorem1.lean}\\
Lemma~\ref{lem:mono} & \texttt{lemmaM}, \texttt{hyp\_of\_le} & \texttt{Monotone.lean}\\
Lemmas~\ref{lem:ore}, \ref{lem:closure} & \texttt{traceable\_of\_traceable\_sup\_edge},\newline
  \texttt{isClosedGraph\_of\_maximal\_nontraceable} & \texttt{Closure.lean}\\
Theorem~\ref{thm:main} & \texttt{conjecture213}, \texttt{conjecture213Statement} & \texttt{Main.lean}\\
\bottomrule
\end{tabular}
\caption{Statements of the paper and their formal counterparts.}\label{tab:lean}
\end{table}

A formal proof is only as good as its statement. The audit file evaluates the definitions on
small graphs, with checks run by the kernel. $\gamma_2(K_{1,3})=3$ and $\gamma_2(P_3)=2$, which
rules out double domination. The lower median of the degree sequence $0,0,1,1$ is $0$, which
rules out the upper median.

\emph{Computer checks.} The proofs above use no computation. As a guard against errors in the
proofs and in the reading of the definitions, their finite consequences were also checked by
computer. Each count below, except those of the random tests in the last item, was obtained by at
least two programs that share no code; the
programs and their outputs are in the supplement to this paper (the ancillary files of its arXiv
version).
\begin{itemize}
\item Theorem~\ref{thm:key} for all closed, connected, non-complete graphs on $n\le 8$ labelled
  vertices. There are none for $n\le 3$, and 4, 35, 1401, 87258 and 9007520 for $n=4,\dots,8$.
  Up to isomorphism, for $n\le 10$, there are 1, 2, 11, 49, 514, 7333 and 209149 for
  $n=4,\dots,10$. For every one of these graphs at least two programs confirm
  $\gamma_2(H)\le\lm(H)+1$, by computing $\gamma_2(H)$ exactly or by checking the set built in
  the proof against the definition. One program checked that set for every choice the proof
  allows.
\item Theorem~\ref{thm:main} for all connected graphs on $n\le 10$ vertices up to isomorphism:
  the hypothesis graphs, counted in the introduction, all have a Hamiltonian path. At $n=10$
  there are 11716571 connected graphs. The counts for $n=9$ and $n=10$ agree with those of the
  independent census of~\cite{aiv}. For $n\le 7$ labelled vertices there are 1, 4, 34, 306,
  13438 and 483928 hypothesis graphs for $n=2,\dots,7$, all traceable.
\item Lemma~\ref{lem:select} for $r\le 6$ by exhaustive search over all systems of $|Y'|$
  subsets of a set $P$ with at most 7 elements, each of size at most $r-2$, in the ranges listed
  in the supplement; for example, there are 61474519 such systems, counted as multisets, with
  $r=|P|=|Y'|=6$.
\item On the small graphs above, the loop of Remark~\ref{rem:procedure} runs at most one round,
  and only on graphs with 10 vertices. It was also run on random closed graphs with up to 64
  vertices, generated with a structure meant to make it run; there it ran for up to 13 rounds. No
  failure occurred.
\end{itemize}

\section*{Sources and status}

We could not reach the catalogue's web site on 21 November 2026. In the Internet Archive's most
recent capture of the list of open conjectures~\cite{wowii}, of 23 July 2026, entry 213 has the
statement above and the status marker~O, as in the captures of 17 July 2010 and 9 November 2016;
the newest capture of the full list, of 26 July 2026, also marks it~O. We read the
definitions in three captures of the catalogue's definitions file (2010, 2024 and 2026), which
agree on the four definitions that entry 213 uses. On 21 November 2026 the Formal Conjectures repository~\cite{fc} had no file, issue or pull
request about entry 213, and a search of arXiv titles and abstracts found nothing about it; neither
did searches of OpenAlex and of Sourcegraph's public code index. We read
the AI Village repository~\cite{aiv} at its last commit, of 6 November 2026 (UTC$-$8). On 21 and
22 November 2026 we searched Google Scholar for ``Conjecture 213'' together with ``Graffiti'' or
``Written on the Wall'', for ``Hamiltonian path'' with ``lower median'' and ``complement'' (from
2024), for ``Graffiti.pc'' with ``2-domination'' and ``Hamiltonian'', and for ``lower median'' with
``2-domination'' and ``traceable''; we also went through the titles and summaries of all results
from 2025 on that mention ``Written on the Wall II''. On 22 November 2026 we searched GitHub's code search for \texttt{GraphConjecture213},
\texttt{wowii-213}, ``WOWII'' with 213, ``Written on the Wall'' with ``Conjecture 213'', and
``Written on the Wall II'' with ``2-domination''. None of these searches found anything about entry
213 beyond the catalogue itself. We did not search GitLab's code search, which requires a login.

\begin{thebibliography}{10}
\bibitem{aiv} AI Village, \emph{WOW-II entry 213: a structural reduction, and the order-10
census}, working note dated 21 October 2026, file
\path{verify/wip/notes_wow2_213_reduction_2026-10-21.md} of the repository
\url{https://gitlab.com/ai-village-agents/village/graffiti-verification}, read at commit
\texttt{8bd9f691fd56} of 6 November 2026 (UTC$-$8).

\bibitem{bc} J.~A. Bondy and V.~Chv\'atal, A method in graph theory, \emph{Discrete Math.}
\textbf{15} (1976) 111--135. \url{https://doi.org/10.1016/0012-365X(76)90078-9}.

\bibitem{gpc} E.~DeLaVi\~na, Graffiti.pc: a variant of Graffiti, in S.~Fajtlowicz, P.~W. Fowler,
P.~Hansen, M.~F. Janowitz and F.~S. Roberts (eds.), \emph{Graphs and Discovery}, DIMACS Ser.
Discrete Math. Theoret. Comput. Sci.~\textbf{69}, Amer. Math. Soc., Providence, RI, 2005, 71--79.
\url{https://doi.org/10.1090/dimacs/069/05}.

\bibitem{wowii} E.~DeLaVi\~na, \emph{Written on the Wall II: Conjectures of Graffiti.pc}, entry
213, dated 12 January 2006, and the definitions file \texttt{wowIIdefs.js}.
\url{http://cms.uhd.edu/faculty/delavinae/research/wowII/}. Consulted on 21 November 2026 in the
Internet Archive capture
\url{https://web.archive.org/web/20260723211215/http://cms.uhd.edu/faculty/delavinae/research/wowII/open.html},
and the capture of the full list
\url{https://web.archive.org/web/20260726061534/http://cms.uhd.edu/faculty/delavinae/research/wowII/all.html}.

\bibitem{dlpw} E.~DeLaVi\~na, C.~E. Larson, R.~Pepper and B.~Waller, Graffiti.pc on the
2-domination number of a graph, \emph{Congr. Numer.} \textbf{203} (2010) 15--32. Read in the
version at \url{https://www.uhd.edu/documents/academics/sciences/gpcon2-dominmay17.pdf}.

\bibitem{fj} J.~F. Fink and M.~S. Jacobson, $n$-domination in graphs, in Y.~Alavi, G.~Chartrand,
L.~Lesniak, D.~R. Lick and C.~E. Wall (eds.), \emph{Graph Theory with Applications to Algorithms
and Computer Science}, Wiley, New York, 1985, 283--300.

\bibitem{fc} Google DeepMind, \emph{Formal Conjectures} repository,
\url{https://github.com/google-deepmind/formal-conjectures}, consulted 21 November 2026.

\bibitem{hp} A.~Hansberg and R.~Pepper, On $k$-domination and $j$-independence in graphs,
\emph{Discrete Appl. Math.} \textbf{161} (2013) 1472--1480. \url{https://doi.org/10.1016/j.dam.2013.02.008}.

\bibitem{hv} A.~Hansberg and L.~Volkmann, Multiple domination, in T.~W. Haynes, S.~T. Hedetniemi
and M.~A. Henning (eds.), \emph{Topics in Domination in Graphs}, Dev. Math.~\textbf{64}, Springer,
Cham, 2020, 151--203. \url{https://doi.org/10.1007/978-3-030-51117-3_6}.

\bibitem{li} R.~Li, Domination number and traceability of graphs, \emph{Discrete Math. Lett.}
\textbf{4} (2020) 27--30. \url{https://www.dmlett.com/archive/v4/DML20_v4_p27_30.pdf}.

\bibitem{lean} L.~de Moura and S.~Ullrich, The Lean 4 theorem prover and programming language,
in \emph{Automated Deduction -- CADE 28}, Lecture Notes in Comput. Sci.~\textbf{12699},
Springer, 2021, 625--635. \url{https://doi.org/10.1007/978-3-030-79876-5_37}.

\bibitem{mathlib} The mathlib Community, The Lean mathematical library, in \emph{Proceedings of
the 9th ACM SIGPLAN International Conference on Certified Programs and Proofs (CPP 2020)}, ACM,
2020, 367--381. \url{https://doi.org/10.1145/3372885.3373824}.

\bibitem{ore} O.~Ore, Note on Hamilton circuits, \emph{Amer. Math. Monthly} \textbf{67} (1960)
55. \url{https://doi.org/10.2307/2308928}.

\bibitem{pepper} R.~Pepper, Implications of some observations about the $k$-domination number,
\emph{Congr. Numer.} \textbf{206} (2010) 65--71.

\bibitem{qing} J.~Qing, The 2-domination number and the upper median degree: a proof of
Graffiti.pc Conjecture 387, arXiv:2607.27246, 28 July 2026.
\url{https://doi.org/10.48550/arXiv.2607.27246}.
\end{thebibliography}

\appendix
\section{The formal statement}\label{app:lean}

The definitions in \texttt{LeanWowii/W213/Defs.lean}, for a finite vertex type \texttt{V} with
decidable equality and a graph \texttt{G : SimpleGraph V} with decidable adjacency, are as
follows; the symbol \texttt{G}$^{\mathsf c}$ is Mathlib's complement.
\begin{alltt}\small
def IsTwoDominating (D : Finset V) : Prop :=
  \(\forall\) v, v \(\notin\) D \(\to\) 2 \(\le\) (D.filter (G.Adj v)).card
theorem exists_twoDominating_card :
  \(\exists\) k, \(\exists\) D : Finset V, IsTwoDominating G D \(\wedge\) D.card = k
noncomputable def gamma2 : \(\mathbb{N}\) :=      -- the least such k
  @Nat.find _ (Classical.decPred _) (exists_twoDominating_card G)
def sortedValues (f : V \(\to\) \(\mathbb{N}\)) : List \(\mathbb{N}\) := (Finset.univ.val.map f).sort (\(\cdot\) \(\le\) \(\cdot\))
def lowerMedianOf (f : V \(\to\) \(\mathbb{N}\)) : \(\mathbb{N}\) :=
  (sortedValues f).getD ((Fintype.card V + 1) / 2 - 1) 0
def lowerMedian (G : SimpleGraph V) [DecidableRel G.Adj] : \(\mathbb{N}\) :=
  lowerMedianOf (fun v => G.degree v)
def Traceable : Prop := \(\exists\) a b : V, \(\exists\) p : G.Walk a b, p.IsHamiltonian

def Conjecture213Statement : Prop :=
  \(\forall\) (V : Type u) [Fintype V] [DecidableEq V] (G : SimpleGraph V) [DecidableRel G.Adj],
    G.Connected \(\to\) 1 < Fintype.card V \(\to\)
    lowerMedian G\textsuperscript{c} + 2 \(\le\) gamma2 G\textsuperscript{c} \(\to\) Traceable G

theorem conjecture213Statement : Conjecture213Statement.\{u\}
\end{alltt}
The index \texttt{(n + 1) / 2 - 1}, with Lean's integer division, is
$\flr{(n+1)/2}-1=\ceil{n/2}-1$; counted from $0$ in a list sorted in nondecreasing order, it
picks the $\ceil{n/2}$-th smallest value. Mathlib's \texttt{Walk.IsHamiltonian} says that every vertex
occurs exactly once on the walk, so a Hamiltonian walk is a Hamiltonian path.

\end{document}
